\emph{પૂર્ણતાનું પ્રમેય} યથાર્થતાના પ્રમેયનું વ્યત્યસ્ત છે. એક સ્વરૂપમાં
$\Gamma \Entails !A$ હોય તો $\Gamma \Proves !A$. બીજા સ્વરૂપમાં,
$\Gamma$ સુસંગત હોય તો તે સંતોષ્ય છે. આપણે બીજું સ્વરૂપ સાબિત કર્યું અને
તેમાંથી પહેલું મેળવ્યું. સાબિતીમાં અનેક પગલાં છે. શરૂઆત સુસંગત ગણ
$\Gamma$થી કરીએ. તેમાં અનંત નવાં અચળચિહ્નો $c_i$ અને
$\lexists[x][!A(x)] \lif !A(c)$ સ્વરૂપનાં સૂત્રો ઉમેરીએ. વિસ્તૃત ભાષામાં
મુક્ત ચલવાળા દરેક સૂત્ર $!A(x)$ને નવા અચળચિહ્નોમાંથી એક સાથે જોડીએ.
આમ $\Gamma$ને સમાવતો \emph{સંતૃપ્ત} અને હજુ સુસંગત એવો વાક્યોનો ગણ મળે
છે. પછી તેને વિસ્તારીને \emph{પૂર્ણ સુસંગત ગણ} બનાવીએ છીએ. આવા ગણમાં
કોઈ પણ વાક્ય $!A$ માટે $!A$ અને $\lnot !A$માંથી બરાબર એક હોય છે. મૂળ
ગણ સંતૃપ્ત હોવાથી હવે આપણી પાસે $\Gamma$ને સમાવતો વાક્યોનો સંતૃપ્ત, પૂર્ણ
અને સુસંગત ગણ $\Gamma^*$ છે. તેના પરથી એવી સંરચના $\Struct{M}$ વ્યાખ્યાયિત
કરી શકાય છે કે $\Sat{M(\Gamma^*)}{!A}$ ત્યારે અને માત્ર ત્યારે જ જ્યારે
$!A \in \Gamma^*$. ખાસ કરીને $\Sat{M(\Gamma^*)}{\Gamma}$, એટલે કે
$\Gamma$ સંતોષ્ય છે. ભાષામાં $=$ હોય તો રચના થોડી વધુ જટિલ બને છે.

પૂર્ણતાના પ્રમેયમાંથી બે અગત્યના ઉપસિદ્ધાંતો મળે છે. \emph{સંહતતાનું
પ્રમેય} કહે છે કે $\Gamma \Entails !A$ ત્યારે અને માત્ર ત્યારે જ જ્યારે
કોઈ સાન્ત $\Gamma_0 \subseteq \Gamma$ માટે $\Gamma_0 \Entails !A$.
સમકક્ષ રીતે, $\Gamma$ સંતોષ્ય છે ત્યારે અને માત્ર ત્યારે જ જ્યારે તેનો
દરેક સાન્ત ઉપગણ $\Gamma_0$ સંતોષ્ય હોય. ચોક્કસ ગુણધર્મો ધરાવતી સંરચનાઓનું
અસ્તિત્વ સાબિત કરવા આ પ્રમેય ઉપયોગી છે. દાખલા તરીકે, મનફાવડા મોટા સાન્ત
નિદર્શો ધરાવતા દરેક સિદ્ધાંતનો અનંત નિદર્શ છે. તેથી પ્રથમ-ક્રમ તર્કમાં
સાન્તતા વ્યક્ત કરી શકાતી નથી. બીજો ઉપસિદ્ધાંત, \emph{લૉવેનહાઇમ--સ્કોલેમનું
પ્રમેય}, કહે છે કે દરેક સંતોષ્ય $\Gamma$નો ગણનીય નિદર્શ છે. તેનાથી જણાય
છે કે પ્રથમ-ક્રમ તર્કમાં અગણનીયતા વ્યક્ત કરી શકાતી નથી.
